\documentclass[11pt]{article}
\usepackage[margin=1in]{geometry}
\usepackage[T1]{fontenc}
\usepackage[utf8]{inputenc}
\usepackage{booktabs}
\usepackage{url}
\title{ES-FA: Event-Driven SNN Acceleration with Local Software-Hardware Evidence}
\author{BTP Progress Draft}
\date{\today}
\begin{document}
\maketitle

\begin{abstract}
This draft summarizes the current ES-FA implementation after the latest local
smoke run. The software pipeline ran successfully with baseline training,
hardware-aware loss, adaptive dataflow, analysis, evidence tables, and output
section generation. The best current software result is
\texttt{exp1\_hardware\_aware\_loss}, which keeps 95.70\% validation accuracy
while reducing the energy proxy by 79.68\% compared with the dense baseline
estimate. Fresh hardware/Vivado execution was skipped in this run; existing
local hardware-validation rows remain available for reference, while KV260 board
measurements are still pending.
\end{abstract}

\section{Motivation}
Edge FPGA inference needs models that are not only accurate but also efficient.
Spiking neural networks are useful for this direction because sparse spike
activity can potentially reduce memory traffic and compute operations. ES-FA
uses a small SNN and tracks both ML metrics and hardware-facing proxy metrics.

\section{Current Implementation}
The repository is organized around the renamed project layout:
\begin{itemize}
\item \texttt{p1\_training/}: SNN model, training core, baseline runner, exports.
\item \texttt{p2\_hardware\_model/}: spike, memory, energy, and latency estimators.
\item \texttt{experiments/}: hardware-aware and adaptive dataflow experiments.
\item \texttt{analysis/}: ranking, plots, and evidence generation.
\item \texttt{hardware/}: Verilog RTL blocks and testbenches.
\item \texttt{hardware\_validation/}: KV260-oriented xsim/Vivado validation flow.
\item \texttt{deployment/} and \texttt{system\_layer/}: CLI and adaptive routing.
\item \texttt{output/}: generated report and presentation material.
\end{itemize}

\section{Model and Estimator}
The current SNN is a compact feed-forward model:
\[
784 \rightarrow 128 \rightarrow 64 \rightarrow 10
\]
with LIF hidden neurons and a default temporal window of 16 time steps. During
training, ES-FA records spike activity, sparsity, memory-access estimates, energy
proxy, and latency proxy.

\section{Latest Smoke-Run Results}
\begin{table}[h]
\centering
\begin{tabular}{lrrrr}
\toprule
Run & Accuracy & Sparsity & Energy proxy & Score \\
\midrule
\texttt{exp1\_hardware\_aware\_loss} & 0.9570 & 0.5648 & 4,592,863.17 & 0.7699 \\
\texttt{baseline\_paper1} & 0.9570 & 0.5648 & 22,604,295.76 & 0.7387 \\
\texttt{exp5\_dataflow\_adaptation} & 0.9570 & 0.5648 & 45,016,894.73 & 0.6699 \\
\bottomrule
\end{tabular}
\caption{Current software ranking from the latest local smoke run.}
\end{table}

The current best run is \texttt{exp1\_hardware\_aware\_loss}. It matches the
baseline accuracy and sparsity in this smoke run, but the energy proxy is much
lower. The proxy reduction compared with the dense baseline is:
\[
\frac{4{,}592{,}863.17 - 22{,}604{,}295.76}{22{,}604{,}295.76}
= -79.68\%.
\]

\section{Hardware Status}
The hardware side contains RTL blocks and a KV260-centered validation flow. The
latest smoke run used \texttt{-SkipVivado -SkipHardware}, so no fresh hardware
build was executed during this run. Existing local evidence rows still include
simulation/report data such as 5760 ns active-window latency and 576 cycles, but
these should be described as previous local validation evidence, not a new board
measurement.

\section{Presentation Claim Boundary}
The safe claim for the current presentation is:
\begin{quote}
The local software and analysis pipeline is working end-to-end. The best
hardware-aware software experiment preserves 95.70\% accuracy and reduces the
energy proxy by 79.68\% against the dense baseline estimate. Hardware RTL and
validation scripts are present, but fresh Vivado hardware execution and KV260
physical measurements remain pending for final hardware claims.
\end{quote}

\section{Next Work}
\begin{itemize}
\item Run xsim-only validation after the presentation smoke run.
\item Run full Vivado validation when build time is available.
\item Collect physical KV260 board measurements.
\item Calibrate estimator constants against measured hardware values.
\item Expand the model/scheduler validation matrix.
\end{itemize}

\section{Conclusion}
ES-FA is now presentation-ready as a local software-hardware research pipeline.
The project can demonstrate training, hardware-aware metrics, analysis plots,
evidence tables, and CLI routing. The remaining work is mainly hardware closure:
fresh Vivado runs, KV260 board measurements, and estimator calibration.

\end{document}
